
\begin{document}
\raggedright
\thispagestyle{fancy}
\begin{center}
{\fontsize{18}{22}\selectfont\bfseries \papertitle\par}
\vspace{11pt}
{\fontsize{14}{17}\selectfont \papertrack\par \paperid\par}
\end{center}

% Nate's abstract (paper/abstract.tex): text verbatim, one paragraph in the Alcorn (2018) format.
% Changes: 5 heads -> 6; section headings, Table/Figure pointers and code link dropped (the floats
% moved to 3.2 and 3.3, the link is in 2.2).
\begin{center}{\fontsize{16}{19}\selectfont\bfseries Abstract}\end{center}
\vspace{-11pt}
In baseball, understanding how to use great hitters is often a simple task: play as many of them as you can and play them often. Teams try to gain an edge with the remaining hitters, but less playing time means thinner data in an area that takes hundreds of plate appearances to stabilize. This paper asks whether a model trained on pitches can project how a hitter will do against a type of matchup, like handedness or a specific pitcher, with less data than standard projections need. We train a neural network inspired by (batter$|$pitcher)2vec that learns player embeddings from hitter-controlled pitch outcomes, which decreases the impact of outcome luck. A plate appearance is broken into 6 heads: swing decision, contact, and quality decomposed as exit velocity, launch angle, and spray direction, each conditioned on the pitcher and the game state. The model trains on 6.5 million Statcast pitches from 2015 through 2024 covering 1,672 hitters, each represented as a 64-dimensional embedding. To query it, we reweight which pitchers the projection is drawn over and compose the outputs through the count-state Markov chain into a projected wOBA, so a new split needs a new weight vector, not a retrained model. We test side-specific wOBA against left- and right-handed pitchers on sealed 2025 data against gradient-boosted and empirical Bayes baselines, using RMSE and weighted rank correlation, grouped by prior plate appearances (PA). The model did not pass our pre-registered test. Pooled RMSE was tied (0.0452 against 0.0453 for gradient boosting and 0.0460 for empirical Bayes) and pooled rank correlation was higher (0.47 against 0.42 and 0.40). It beat both baselines on both measures only for hitters with 113 to 451 prior plate appearances, an exploratory result. On 2024, the season used to select the model, it recovers 60\% of the rank-order signal and 64\% of the level signal that one season of data can measure. Without being given definitions of power or contact, the embedding still separates hitters on both, with held-out accuracy of 0.79 and 0.75 against 0.33 by chance, and observed wOBA rises toward high power and high contact. The model is better at ordering hitters within a split than at predicting how many points of wOBA the split is worth. For a team, that makes it most useful for ranking platoon options among hitters with a moderate track record, where the decision is actually in question, not as a replacement for a projection system on raw accuracy.

% =====================================================================================
\section{Introduction}
\label{sec:intro}

% Source: src/analysis/claim1_eval.py STABILIZATION_N_STAR = 226 (variance-components estimate,
% training seasons, hitters only).
Most lineup decisions make themselves. A team's best hitters play every day against every kind of pitcher. The decisions that leave room for improvement involve the rest of the roster, the bench bats and platoon hitters, and they usually turn on a split, meaning how a hitter performs against one kind of pitcher. These hitters also have the least data. In the training seasons, a hitter's wOBA against left-handed pitchers takes about 226 plate appearances (PA) against them before half of its variation across hitters reflects skill rather than noise (Section~\ref{sec:ceiling}). A part-time hitter can take several seasons to get there.

Much of that noise enters at the end of the plate appearance. Fielders, parks, and where a ball in play happens to land all move a wOBA outcome without reflecting anything the hitter did. The hitter's own contribution comes earlier and more often. On each pitch the hitter decides whether to swing and, after swinging, either makes contact or misses, and Statcast records the exit velocity, launch angle, and direction of nearly every ball in play. A plate appearance yields one wOBA value but several of these pitch-level outcomes, and they carry less of the luck that follows contact.

% Source: alcorn-2018-batter-pitcher2vec.pdf Sections 2.1, 2.2, 4; tango-2009 platoon constants note
% (The Book p.157).
\citet{alcorn2018} showed that a neural network can learn useful representations of players from their matchups alone. The (batter$|$pitcher)2vec model gave each batter and pitcher a nine-dimension embedding, trained to predict the outcome of each at-bat from the identities of the two players in Retrosheet data from 2013 through 2015. The embeddings grouped similar hitters without being told what made them similar, and they predicted outcome distributions for matchups absent from training better than simpler baselines. Alcorn suggested that predicting PITCHf/x measurements instead of discrete outcomes could convey more about players. The model described here moves in that direction with Statcast. Each hitter keeps a learned embedding, but the network predicts the hitter's response to each pitch from the pitch's measured characteristics, rather than the at-bat result from the pitcher's identity.

The usual alternative for a split is to regress the observed split toward the league split for the hitter's side. In \citet[p.~157]{tango2007}, a right-handed hitter needs about 2,200 PA against left-handed pitchers before the observed split receives half the weight, and a left-handed hitter about 1,000, so for most hitters the estimate stays close to the league-average split. Empirical Bayes shrinkage of this kind \citep{efron1972} is one of the two baselines in this paper.

This paper introduces a pitch-level hitter embedding model and a query that turns it into a projected wOBA against any chosen set of pitchers. Unlike split regression, which uses only a hitter's outcomes in that split, the model learns from every pitch the hitter has seen and shares structure across hitters through the embedding, so a new split requires a new set of pitcher weights rather than a retrained model. The query is tested on 2025 handedness splits against empirical Bayes and gradient-boosted baselines, and the embedding is examined for what it learned from per-pitch outcomes alone.

% =====================================================================================
\section{Methods}
\label{sec:methods}

\subsection{Data}
\label{sec:data}

% Source: data/processed/phase_d5_clean_final/manifest.json; src/config/split_config*.json;
% decision log 2026-08-01, 2026-09-20.
Pitch-level Statcast data from the 2015 through 2025 regular seasons were used. Each pitch carries its velocity, movement, release point, spin, plate location, pitch type, the count, and the handedness of both players, and each ball in play carries its exit velocity, launch angle, and a spray angle computed from where it was fielded. Pitchouts, automatic balls and strikes, bunts, pitches thrown by position players, and pitches missing a core measurement were removed, as were pitchers batting. No minimum exposure was imposed, so every hitter in the training seasons receives an embedding, and hitters absent from them share one reserved embedding. The cleaned table holds 7,245,923 pitches. Model configurations were compared by training on 2015 through 2023 and scoring on 2024. The chosen configuration was then retrained on 2015 through 2024, 6,540,616 pitches covering 1,672 hitters and 2,182 pitchers, and evaluated on 2025. Feature scaling was fit on 2015 through 2023 only, and baselines see only seasons before the one being scored.

\subsection{Model}
\label{sec:model}

% Source: paper/draft3/fig_arch.tex
\begin{figure}[!ht]
\centering
% fig_arch.tex -- architecture (a) and query (b) schematic. \input inside a figure env.
% needs: tikz with arrows.meta, positioning, fit, calc, shapes.geometric, backgrounds; amsmath; xcolor
\begin{tikzpicture}[
  x=1cm, y=1cm, font=\scriptsize, text=black!80, line width=0.5pt,
  box/.style={draw=black!55, rounded corners=1.5pt, align=center, inner sep=2pt,
              minimum height=0.5cm},
  acc/.style={draw=blue!60, fill=blue!7},
  bar/.style={draw=black!55, fill=black!5, rounded corners=1.5pt,
              minimum width=0.46cm, minimum height=5.6cm},
  head/.style={box, minimum width=3.2cm, anchor=west},
  arr/.style={-{Stealth[length=3.5pt,width=3pt]}, draw=black!60},
  cond/.style={-{Stealth[length=3pt,width=2.6pt]}, draw=blue!60, densely dashed},
  ttl/.style={anchor=west, font=\footnotesize\bfseries},
  note/.style={text=black!60},
  cnt/.style={draw=black!45, rounded corners=1pt, inner sep=0pt, minimum width=0.42cm,
              minimum height=0.26cm, font=\tiny},
  stepc/.style={circle, draw=blue!60, fill=blue!7, inner sep=0pt, minimum size=0.38cm,
                font=\scriptsize\bfseries},
  para/.style={anchor=north west, align=left, inner sep=0pt},
]

% ---------------------------------------------------------------- (a) the network
\node[ttl] at (0,8.55) {(a) The network};

\node[box, acc, minimum width=2.45cm] (hid) at (1.25,7.55) {Hitter ID $h$};
\node[box, acc, minimum width=2.45cm] (emb) at (1.25,6.65) {Embedding (64)};
\node[box, minimum width=2.45cm, text width=2.3cm] (ctx) at (1.25,4.9)
  {Pitch context (46)\\[1pt]{\tiny\color{black!60}physical measurements,
   location, type, handedness of both players, count}};
\node[box, minimum width=2.45cm] (c1) at (1.25,3.45) {Context layer (128)};
\node[box, minimum width=2.45cm] (c2) at (1.25,2.55) {Context layer (128)};
\draw[arr] (hid) -- (emb);
\draw[arr] (ctx) -- (c1);
\draw[arr] (c1) -- (c2);

\node[bar] (cat) at (3.05,5.0) {};
\node[bar] (t1)  at (3.8,5.0)  {};
\node[bar] (t2)  at (4.55,5.0) {};
\node[rotate=90] at (cat) {Concatenate (64 + 128)};
\node[rotate=90] at (t1)  {Trunk layer (256)};
\node[rotate=90] at (t2)  {Trunk layer (256)};
\draw[arr] (emb.east) -- (cat.west |- emb.east);
\draw[arr] (c2.east)  -- (cat.west |- c2.east);
\draw[arr] (cat) -- (t1);
\draw[arr] (t1) -- (t2);

% six heads, top to bottom in plate-appearance order
\foreach \name/\y/\lbl in {
    sw/7.45/{Swing},
    co/6.47/{Contact $\mid$ swing},
    sp/5.49/{Foul/tip/in play $\mid$ contact},
    ev/4.51/{Exit velocity {\color{black!60}(24 bins)}},
    la/3.53/{Launch angle $\mid$ EV {\color{black!60}(24)}},
    sy/2.55/{Spray $\mid$ EV, LA {\color{black!60}(24)}}} {
  \node[head] (\name) at (5.1,\y) {\lbl};
  \draw[arr] (t2.east |- \name.west) -- (\name.west);
}
% EV and LA bins feed the later heads
\draw[cond] (ev.east) to[out=0, in=0, looseness=1.3] (la.east);
\draw[cond] (ev.east) to[out=0, in=0, looseness=1.1] (sy.east);
\draw[cond] (la.east) to[out=0, in=0, looseness=1.3] (sy.east);

\node[box, dashed, draw=black!50, text width=8.1cm, inner sep=3pt] (take) at (4.25,1.3)
  {Take table, fit separately and not part of the network\\[1pt]
   $p(\text{ball},\ \text{called strike},\ \text{HBP} \mid
     \text{location, batter hand, pitcher hand})$};

% panel divider
\draw[black!20] (8.85,0.8) -- (8.85,8.7);

% ---------------------------------------------------------------- (b) a query
\begin{scope}[xshift=9.15cm]
\node[ttl] at (0,8.55) {(b) A query};

% step 1
\node[stepc] (s1) at (0.2,7.9) {1};
\node[para, text width=6.4cm] at (0.55,8.03)
  {For pitcher $j$ and each count, sample 6 of the pitcher's real pitches in that count
   to same-side hitters. Run the network for hitter $h$ and average.\\[3pt]
   {\color{blue!60}$\Rightarrow$} per count $p_B$, $p_S$, $p_F$, $p_H$ and $v(b,s)$,
   the expected wOBA from a ball in play\\(its chance times its average value)};

% step 2, the 12-count chain
\node[stepc] (s2) at (0.2,5.5) {2};
\begin{scope}[shift={(0.95,5.5)}]
  \foreach \b in {0,...,3} \foreach \s in {0,1,2}
    \node[cnt] (n\b\s) at (0.62*\b,-0.48*\s) {\b-\s};
  \foreach \b in {0,1,2} \foreach \s in {0,1,2} {
    \pgfmathtruncatemacro{\nb}{\b+1}
    \draw[arr, draw=black!40] (n\b\s) -- (n\nb\s);}
  \foreach \b in {0,...,3} \foreach \s in {0,1} {
    \pgfmathtruncatemacro{\ns}{\s+1}
    \draw[arr, draw=black!40] (n\b\s) -- (n\b\ns);}
  \node[cnt, acc] at (n00) {0-0};
  \node[note, font=\tiny, anchor=south] at ($(n10.north)!0.5!(n20.north)+(0,0.03)$)
    {balls $\rightarrow$};
  \node[note, font=\tiny, anchor=north] at (n12.south |- 0,-1.15) {strikes $\downarrow$};
\end{scope}
\node[para, text width=3.45cm] at (3.3,5.63)
  {Count-state Markov chain over the 12 counts, solved backward from the end states\\[3pt]
   {\color{blue!60}$\Rightarrow$} $W_{hj} = W(0,0)$};

% step 3
\node[stepc] (s3) at (0.2,3.45) {3};
\node[para, text width=6.4cm] at (0.55,3.58)
  {Weighted mean over the pitcher set $S$, with weights proportional to batters faced};
\node[anchor=west] at (0.5,2.42)
  {$\widehat{\text{wOBA}}_S(h)=\sum_j w^S_j\, W_{hj}, \qquad
    w^S_j\propto \text{BF}_j\,\mathbf{1}[j\in S]$};
\node[box, draw=blue!60, text width=6.2cm, inner sep=3pt, anchor=north west]
  at (0.55,1.98) {A new split is a new set $S$, so only the weights change.
                 No retraining.};

\draw[arr] (s1) -- (s2);
\draw[arr] (s2) -- (s3);
\end{scope}
\end{tikzpicture}

\caption{\toreview (a) The network. A hitter's embedding and the pitch context feed a shared trunk, and six heads predict the hitter-controlled outcomes of the pitch in the order they occur. Taken pitches are read from a separate table. (b) A query. Per-count outcomes against one pitcher are composed through the count into $W_{hj}$, and a weighted average over the chosen pitchers gives the projection. A new split changes only the weights.}
\label{fig:arch}
\end{figure}

% Source: src/model/v1.py heads; data manifest n_quality_bins 24; decision log 2026-08-08.
The network is shown in Figure~\ref{fig:arch}a. Each hitter has a 64-dimension embedding, a learned vector that is the only input identifying the hitter. Each pitch is described by a 46-dimension context vector built from its physical measurements, location, type, the handedness of both players, and the count. The pitcher enters only through the pitches thrown, not through an identifier, so the model can describe any pitcher whose pitches it has seen. The count is the only game state included, so predictions reflect the hitter and pitch rather than the game situation.

The context passes through two layers of 128 units and is joined with the hitter's embedding in a shared trunk of two layers of 256 units. Six heads read the trunk in the order a plate appearance unfolds. They give the probability of a swing, the probability of contact given a swing, whether contact is a foul, a foul tip, or a ball in play (a two-strike foul tip is a strikeout and a foul is not), and the exit velocity, launch angle, and spray direction of a ball in play, each in 24 bins. Launch angle is predicted given exit velocity, and spray given both, so the three are never treated as independent. Taken pitches depend mostly on the pitcher and the umpire, so their outcomes come from a table of the probability of a ball, a called strike, or a hit batter by plate location and by the pairing of batter and pitcher hand, fit on the training seasons.


% Source: src/model/v1.py loss; src/model/train.py masks.
The network is trained to maximize the summed log-likelihood of each hitter's outcomes,
\begin{equation}
\mathcal{L} = \sum_{i} \Big[ \log p(a_i) + m^{c}_i \log p(c_i) + m^{t}_i \log p(t_i \mid c_i) + m^{e}_i \log p(e_i) + m^{l}_i \log p(l_i \mid e_i) + m^{s}_i \log p(s_i \mid e_i, l_i) \Big],
\label{eq:loss}
\end{equation}
where for pitch $i$, $a_i$ is the swing decision, $c_i$ the contact outcome, $t_i$ the foul, foul tip, or ball in play split, and $e_i$, $l_i$, and $s_i$ the exit velocity, launch angle, and spray bins. Every probability is conditioned on the hitter $h$ and the pitch context $x_i$. Each mask $m_i$ is one when the outcome exists and was measured, so contact is scored only on swings, the split only on contact, and the three quality terms only on measured balls in play.

% Source: results/model_v1/sweep_log.csv stage clean (dim64 five-seed mean 1.023268, SD 0.000128;
% base 1.02378; dropout 1.02357); decision log 2026-09-02/03 (SGD table, lr 1.0); checkpoint
% clean_dim64_final_s0.pt (266,317 parameters); review-pack-refit.md (18,538 steps).
The network has 266,317 parameters. Weights outside the embedding table are trained with AdamW \citep{loshchilov2019} at a learning rate of $10^{-3}$ and batches of 8,192 pitches. Adaptive and plain stochastic gradient descent were both tested on the embedding table, because an adaptive optimizer gives a rarely seen hitter steps as large as a regular's. Plain stochastic gradient descent at a learning rate of 1.0 was chosen because it moves each embedding in proportion to the hitter's data, so embedding length grows with exposure, and on an earlier build of the data it gave lower held-out loss. Embedding sizes of 32 and 64 were compared, as the architecture plan called for, and 64 was chosen because it gave the lower 2024 log loss averaged over five seeds (1.02327 against 1.02378, a gap about four times the seed-to-seed standard deviation of the 64-dimension runs). Weight decay of $10^{-2}$ on all weights, including the embedding table, was kept after stronger and weaker settings each scored worse in a single-seed screen, and dropout of 0.1 in the trunk was fixed in advance. The final model was trained for 18,538 steps, the median early-stopping point of the five selection runs, with their learning-rate reductions replayed at the same steps. Five copies are trained from different seeds as a deep ensemble \citep{lakshminarayanan2017}, and each head's predicted distribution is averaged across the five before any query. The five selection runs give similar vectors for the same hitter, with a median cosine similarity of 0.83 after rotating each onto the first, against 0.12 when hitters are shuffled before the rotation is fit. The model was implemented in PyTorch and trained on a CPU in about 25 minutes per seed. The code is available at \url{https://github.com/nseluga/hitter-embedding}.

\subsection{Querying the model}
\label{sec:query}

% Source: src/model/query.py docstring; src/model/query_tables.py; decision log 2026-08-08.
The network describes a single pitch, while a lineup decision concerns many plate appearances against many pitchers. A query connects the two in three steps (Figure~\ref{fig:arch}b).

The first step fixes one pitcher. For pitcher $j$ and each count, six real pitches that pitcher threw in that count to hitters of the same side during the training seasons are run through the network for hitter $h$ and averaged. Sparse counts fall back to pitches with the same number of strikes, then to all of that pitcher's pitches to that side. This gives, for each count, the probability that the pitch ends as a ball, a strike, a foul, a hit batter, or a ball in play, and the expected wOBA from a ball in play, its probability times its average value. The average value sums, over every combination of exit velocity, launch angle, and spray bins, the hitter's predicted probability of that combination times the league's average wOBA for it. Balls in play that Statcast did not measure, about 7\%, take the league's average value for such balls.

The second step plays out the plate appearance. The count moves as a Markov chain over 12 states from 0-0 to 3-2, and the expected wOBA from count $(b,s)$ satisfies
\begin{equation}
W(b,s) = p_B\, W(b+1,s) + p_S\, W(b,s+1) + p_F\, W(b,\min(s+1,2)) + p_H\, w_{\text{HBP}} + v(b,s),
\label{eq:chain}
\end{equation}
where $p_B$, $p_S$, $p_F$, and $p_H$ are the probabilities of a ball, a strike, a foul, and a hit batter in that count, $v(b,s)$ is the expected value from a ball in play, a fourth ball is worth a walk's wOBA weight, and a third strike is worth zero. A two-strike foul leaves the count unchanged, and that loop is solved directly. Working back from the terminal states gives $W_{hj} = W(0,0)$, the projected wOBA of hitter $h$ in one plate appearance against pitcher $j$, with no simulation.

The third step chooses the pitchers. A query is a set of pitchers $S$, and the projection is a weighted average over it,
\begin{equation}
\widehat{\text{wOBA}}_S(h) = \sum_{j} w^S_j \, W_{hj}, \qquad w^S_j \propto \text{BF}_j \, \mathbf{1}[j \in S],
\label{eq:query}
\end{equation}
where $\text{BF}_j$ is pitcher $j$'s batters faced in the training seasons. A platoon split takes $S$ as all left-handed or all right-handed pitchers. Any other pitcher set uses the same calculation with no retraining. Only the handedness query is tested against outcomes here, and Section~\ref{sec:otherqueries} shows what two other queries return.

\subsection{Evaluation}
\label{sec:eval}

% Source: src/analysis/claim1_eval.py:95-108 (n* 226, cuts (113, 452), floor 10 PA);
% min-PA sweep verdict JSON.
The handedness query is scored on a hitter's wOBA against one pitcher hand in 2025, called a unit. The model was built mainly to help with hitters who have little history against a hand, but it is meant to be accurate for every hitter, so results are reported for all units and by prior exposure, the hitter's PA against that hand in all earlier seasons since 2015. Units fall into three strata, low (fewer than 113 prior PA), medium (113 to 451), and high (452 or more). The cuts are half and twice the 226 PA at which half the variation in wOBA against left-handed pitchers reflects skill, a simple convention rather than a tuned choice. A unit needs at least 10 PA in 2025. The same projections also give a hitter's overall level (wOBA pooled over both hands) and platoon split (wOBA against left-handed pitchers minus wOBA against right-handed pitchers).

% Source: src/analysis/baseline_ladder_bivariate_eb.py (3 trailing seasons, covariance by batter type);
% claim1_eval.py; paired CSV (2,000 batter-clustered draws); decision log 2026-07-30, 2026-09-20.
Two baselines are used. The empirical Bayes baseline takes a hitter's wOBA against each hand over the previous three seasons and shrinks the two sides jointly toward the league average for the hitter's batting side, with the correlation between sides estimated from the data. It generalizes the platoon regression of \citet{tango2007} by learning the amount of shrinkage. The boosted baseline is gradient-boosted trees \citep{chen2016} fit on one row per hitter, season, and pitcher hand, using the same trailing outcomes plus summaries of the hitter's own pitches, such as swing, whiff, and exit velocity rates by pitcher hand and pitch group. It does not see individual pitches or pitchers, and the count enters only as a two-strike indicator. Each method is scored by root mean squared error (RMSE) and by Spearman rank correlation, both weighted by each unit's 2025 PA. Differences between the model and a baseline carry 95\% intervals from 2,000 bootstrap draws that resample batters. The low stratum was specified as the primary test before the 2025 model was trained. The model had to beat the boosted baseline on RMSE and both baselines on rank correlation. All other comparisons are exploratory.\footnote{2025 was also scored with an earlier build that used a 32-dimension embedding and kept pitchers batting. No comparison changed whether its interval excluded zero, although the low-stratum error differences moved from nominally favoring the baselines to nominally favoring the model.}

